\begin{abstract}
We study the number of persistent states needed to realize a causal
statistical response. For a complete finite response array we define an
intrinsic support-component rank from all normalized continuations.
Every stochastic realization obeys this cutwise lower bound. If the
convex hull in each component is a simplex, residual vertices attain
all these bounds simultaneously through normalized causal updates.
The criterion is finite and effective on rational data, and the complete
class is closed under independent products of terminal channels, with
multiplicative state counts. A growing binary random-access experiment
then separates predictive dimension from positive memory: its exact
minimum whole-schedule peak is $n+1$ for $0<\eta\le n^{-2}$ and $2^n$
for $1-n^{-1}<\eta\le1$, even when acquisition order is chosen
adaptively. The weak-signal bound is attained by an explicit sequence
of compatible enclosing simplices, not merely by a final encoder.
Minimax contact connects the classification to the marked physical
audit and its exact five-state decision law. The previously established
positive-noise saddle, fixed-clock validation-order classification,
support-boundary jump, causal foundations and transport proofs are
retained as supporting mathematics. Exact physical counts assume a
specified target and atomic stochastic rows; acquired calibration has
separately priced approximation guarantees.
\end{abstract}
\maketitle
\tableofcontents
\clearpage

\section{Introduction}
A finite stochastic experiment can have a low-dimensional predictive
image and nevertheless require many persistent states. The reason is
positivity: an internal state describes a normalized future experiment,
not an arbitrary vector in a linear span. Causality imposes a second
constraint. Positive representations at successive cuts must be joined
by common stochastic updates, and a schedule chosen from observations
cannot reappear as an uncharged clock.

We seek conditions under which those constraints admit an intrinsic
and exact state count. A lower bound obtained from selected faces and
an independently supplied upper program can certify a particular
example. It does not yet tell us when the continuation geometry itself
determines a complete answer. Theorem~\ref{thm:intrinsic29} supplies
such a condition for a finite class of causal experiments. Its data are
all normalized continuation rows of the response array, their minimal
faces in the causal polytope, and the connected components of the
face-intersection graph. The resulting sum of affine ranks is a lower
bound against every stochastic realization. If the residual convex
hull in each component is a simplex, a canonical residual-vertex
construction attains those ranks at every cut in one machine.

The condition does not assume a dimension-growth law, a metric entropy
exponent, or a proposed controller. It can be checked by exact linear
feasibility and rank. It is also not universal: an explicitly given
full-support square has rank three and needs four states. The failure
example makes clear what the completeness hypothesis excludes.
For independent terminal channels the complete class is closed under
products, and its intrinsic capacities multiply
(Theorem~\ref{thm:tensor29}). At every declared schedule vertex this
gives an exact field rather than a field of analyst-selected
certificates. The schedule optimum follows by a bottleneck recurrence.

A second theorem gives a more substantial separation in a growing
experiment. Read $n$ bits, then answer a later coordinate query with
specified bias $\eta$. Every input and every query must have exactly
the required conditional output law. The checkpoint geometry is the
containment of $[-\eta,\eta]^n$ in a convex hull inside $[-1,1]^n$.
For $0<\eta\le n^{-2}$, a sequence of enclosing simplices has
$(t+1)$ vertices after $t$ reports and is closed under appending the
next bit. Its barycentric updates implement the channel with peak
$n+1$, which the affine dimension forces. For $\eta>1-n^{-1}$,
disjoint corner caps force $2^n$ states, attained by full retention.
Both lower bounds hold after any adaptive acquisition schedule,
because the whole acquired transcript must pass through the state
before the query. Both regimes include full-support laws. Thus the
same predictive rank $n+1$ can be attained by causal memory or be
exponentially smaller than it.

The statistical application enters in the opposite direction: a
minimax saddle first forces continuations, and the intrinsic geometry
then prices them. Bayes indifference alone does not specify an optimal
continuation. For the marked two-preparation family, contact and
stationarity fix two coordinates left free by a least-favorable prior.
The five forced rows have canonical support components of dimensions
$1,0,1$, recovering the exact count $2+1+2=5$. These rows remain forced
in the full-support physical rectangle. Their independent query-channel
products have capacity $5^m$; this is a statement about specified
product continuations, not a new formula for the multi-block minimax
value.

The proof structure is therefore
\[
\begin{gathered}
 \text{uniform saddle contact}
 \ \longrightarrow\ \text{forced normalized continuations}\\
 \longrightarrow\ \text{intrinsic positive geometry}
 \ \longrightarrow\ \text{compatible simplex updates}\\
 \longrightarrow\ \text{one causal realization with matching memory}.
\end{gathered}
\]
The general positive-realization and residual-automaton mechanisms have
classical antecedents. Section~\ref{sec:literature29} compares the
precise claims with stochastic realization, residual automata,
nested-polytope factorization, incomplete-machine minimization and
random-access coding. We do not claim a new general minimax interchange
or a solution of unrestricted positive-rank minimization.

\subsection{Resource and parameter conventions}
All persistent randomness belongs to the register. A fixed finite
program and known row constants are read-only; an unknown target is
not a free constant. The finite-state theorems use atomic stochastic
rows, not free finite-length exact sampling of arbitrary real coins.
The random-access query is one atomic input, explicitly distinct from
a bit-serial query buffer. Input-dependent acquisition choices are
allowed, but their retained protocol information is charged. The
physical twelve-state optimum is over externally declared validation
clocks; no adaptive collision optimum is inferred from it.

For a specified positive-noise Brownian target the exact physical saddle
and five/twelve-state conclusions are retained. Acquired target
calibration instead gives an approximate-score controller with finite
counter and precision costs. The main theorems keep these quantifiers
separate.

\subsection{Organization and the retained development}
Sections~\ref{sec:intrinsic29}--\ref{sec:literature29} contain the
current theorem spine and its proofs. The supporting sections then
retain the previous causal foundations, exact saddles, physical
identification and validation-order proofs, revelation laws,
finite-architecture certification, all-preparation localization, and
transport consumers. Their stable theorem labels remain resolvable in
this single manuscript. No predecessor proof is discarded to shorten
the main narrative. The historical A2 geometric root and B4/C2 operator
gates retain their own hypotheses; neither the present class theorem
nor the growing query example replaces those independent derivations.
